\documentclass[11pt,a4paper]{article}

% --- Packages ---
\usepackage[utf8]{inputenc}
\usepackage[margin=1in]{geometry}
\usepackage{microtype}
\usepackage{hyperref}
\usepackage{booktabs}
\usepackage{xcolor}
\usepackage{graphicx}
\usepackage{float}
\usepackage{amsmath,amssymb}
\usepackage{enumitem}
\usepackage{titlesec}
\usepackage{caption}

% --- Colors ---
\definecolor{primary}{RGB}{11, 37, 48}
\definecolor{accent}{RGB}{2, 132, 199}
\definecolor{tableheader}{RGB}{235, 242, 244}
\definecolor{darktext}{RGB}{30, 41, 59}

% --- Hyperref Setup ---
\hypersetup{
    colorlinks=true,
    linkcolor=accent,
    filecolor=accent,
    urlcolor=accent,
    citecolor=accent
}

% --- Section Styling ---
\titleformat{\section}
  {\normalfont\large\bfseries\color{primary}}{\thesection.}{0.6em}{}[\vspace{-0.4em}\color{accent}\rule{\linewidth}{1pt}]

\titleformat{\subsection}
  {\normalfont\normalsize\bfseries\color{primary}}{\thesubsection}{0.6em}{}

\setlist[itemize]{noitemsep, topsep=2pt, leftmargin=1.5em}

% --- Document Info ---
\title{\vspace{-1.5cm}\textbf{\LARGE Spotter Freight Rate Prediction}\\[0.3em]
\Large Technical Assessment Report}
\author{\textbf{Machine Learning Engineer Assessment Candidate}}
\date{September 2025}

\begin{document}

\maketitle
\vspace{-0.8cm}

\section{Problem Overview \& Objectives}
The objective of this assessment is to develop an accurate, robust Machine Learning pipeline to predict spot market truckload freight rates (\texttt{posted\_rate} in USD) from load-level features (origin/destination cities, mileage, equipment type, weight, and date information). 

The primary deliverables include:
\begin{itemize}
    \item \textbf{12,000 Validation Predictions}: Out-of-sample rate predictions for unlabeled loads in \texttt{validation.csv} (\texttt{TE-000001} to \texttt{TE-012000}).
    \item \textbf{31-Day December Forecast}: Daily rate forecasts for a benchmark corridor (Lexington, KY $\rightarrow$ Fort Wayne, IN; 360 miles, Dry Van, 32,000 lb) across December 2025.
    \item \textbf{Codebase \& Scorer Compatibility}: Clean, reproducible training and inference modules conforming to the official \texttt{score.py} validation suite.
\end{itemize}

\section{Dataset Summary \& Quality Findings}
The modeling pipeline is developed using the historical dataset \texttt{train\_test.csv} and evaluated against \texttt{validation.csv}:
\begin{itemize}
    \item \textbf{Historical Labeled Data (\texttt{train\_test.csv})}: 48,000 shipments dated from January 1, 2025 to October 31, 2025 (304 consecutive days), containing 14 attributes including mileage, weight, equipment, coordinates, market signals, and ground-truth \texttt{posted\_rate}.
    \item \textbf{Validation Data (\texttt{validation.csv})}: 12,000 unlabeled shipments spanning November 1, 2025 to December 31, 2025.
    \item \textbf{Target Distribution}: \texttt{posted\_rate} is strictly positive (\$57 to \$25,533) and heavily right-skewed ($\text{skew} \approx 1.90$). Modeling $\log(1 + y)$ was adopted to stabilize gradients and reduce penalty on extreme tail values.
    \item \textbf{Mileage Dominance}: Trip distance is the single strongest rate predictor ($r \approx 0.91$), exhibiting mild non-linearity that benefits from a $\log(1 + \text{distance})$ transformation.
    \item \textbf{Data Cleaning}: Identified 292 negative weight records caused by sign-entry errors (e.g., $-35,000$\,lb). These were resolved via absolute value transformation ($\text{abs}(\text{weight})$). Missing values in \texttt{weight} (0.62\%) and \texttt{market\_index} (0.78\%) were preserved for pipeline median imputation. No training rows were dropped.
\end{itemize}

\section{Feature Engineering \& Leakage Prevention}
\subsection{Feature Transformation Pipeline}
\begin{itemize}
    \item \textbf{Distance}: Log-transformed via $\log(1 + \text{distance})$ followed by standard scaling.
    \item \textbf{Weight}: Sign-corrected, median-imputed, and standard-scaled.
    \item \textbf{Coordinates}: Median-imputed pickup and delivery latitudes/longitudes, standard-scaled.
    \item \textbf{Categoricals}: One-hot encoding on \texttt{equipment}, \texttt{pickup}, and \texttt{delivery} with unknown category handling.
    \item \textbf{Lane Frequency}: Synthesized $\text{lane} = \text{pickup} \rightarrow \text{delivery}$ and frequency-encoded based on historical volume.
    \item \textbf{Cyclical Calendar}: Sine and cosine periodic transformations for month ($T=12$), quarter ($T=4$), and day-of-week ($T=7$) to capture recurring demand cycles.
    \item \textbf{Excluded Columns}: \texttt{load\_id} (identifier), \texttt{market\_index}, and \texttt{quote\_signal} (dropped to maintain pipeline simplicity and prevent noise).
\end{itemize}

\subsection{Preventing Data Leakage}
To ensure unbiased performance estimates, the dataset was strictly split by date \textit{before} fitting any transformers. All scalers, imputers, encoders, and lane frequency mappings were fitted exclusively on the training fold ($\text{date} \le \text{2025-08-31}$). Unseen lanes in holdout or test data default to the minimum frequency observed in training. Validation and December datasets were transformed strictly with the pre-fitted pipeline without refitting.

\section{Validation Strategy \& Model Benchmarks}
\subsection{Time-Based Split Design}
Freight rates exhibit macro market drift and seasonal demand shifts; a standard random K-Fold split would leak future market dynamics into the past. We implemented a chronological split at the 80th percentile of labeled dates:
\begin{itemize}
    \item \textbf{Training Fold}: Dates $\le$ August 31, 2025 (38,477 rows, Jan–Aug 2025).
    \item \textbf{Internal Validation Fold}: Dates $>$ August 31, 2025 (9,523 rows, Sep–Oct 2025).
\end{itemize}
All metrics are reported on the original USD dollar scale after reversing the log transform via $\exp(\hat{y}) - 1$.

\subsection{Model Comparison}
Six model architectures were benchmarked on the internal holdout set under identical preprocessing:

\begin{table}[H]
\centering
\small
\renewcommand{\arraystretch}{1.15}
\begin{tabular}{lcccc}
\toprule
\textbf{Model Architecture} & \textbf{Holdout MAE (\$)} & \textbf{Holdout RMSE (\$)} & \textbf{Holdout $R^2$} & \textbf{Verdict} \\
\midrule
\textbf{CatBoost (Tuned)} & \textbf{133.0854} & \textbf{641.1042} & \textbf{0.8235} & \textbf{Selected Winner} \\
CatBoost (Baseline) & 143.1829 & 644.3404 & 0.8217 & Strong baseline \\
LightGBM & 145.9222 & 644.9609 & 0.8214 & Competitive boosting \\
Linear Regression & 162.1602 & 650.9376 & 0.8181 & Higher linear residual \\
Random Forest & 167.6709 & 662.6681 & 0.8114 & Suboptimal depth \\
XGBoost & 167.6918 & 660.4121 & 0.8127 & Higher tail error \\
\bottomrule
\end{tabular}
\caption{Model performance comparison on the internal time-based holdout set (Sep--Oct 2025).}
\end{table}

\section{Final Model \& Hyperparameter Configuration}
\textbf{CatBoostRegressor} was selected as the final model due to its superior generalization, symmetric tree structure, and built-in categorical regularization. Hyperparameters selected via TimeSeriesSplit grid search:
\begin{itemize}
    \item \texttt{depth}: 6
    \item \texttt{iterations}: 500
    \item \texttt{learning\_rate}: 0.03
    \item \texttt{l2\_leaf\_reg}: 1
    \item \texttt{loss\_function}: 'RMSE'
    \item \texttt{random\_state}: 42
\end{itemize}

\section{December 2025 Prediction Analysis}
The December test corridor evaluates a fixed operational shipment: Lexington, KY to Fort Wayne, IN (360 miles, Dry Van, 32,000 lb). City coordinates were retrieved from the training database. Because route, mileage, load weight, equipment, month ($12$), and quarter ($4$) remain constant across all 31 rows, the model's predicted rates (\$769.92 to \$787.05) cleanly isolate and capture the learned \textbf{day-of-week weekly freight rate cycle}.

\begin{figure}[H]
    \centering
    \includegraphics[width=0.92\linewidth]{candidate_december.png}
    \caption{December 2025 predicted rate trajectory for the Lexington $\rightarrow$ Fort Wayne corridor generated by \texttt{score.py}.}
\end{figure}

\section{Limitations \& Considerations}
\begin{itemize}
    \item \textbf{Validation Ground Truth}: The reported metrics reflect the internal Sep--Oct 2025 holdout; Spotter evaluates the true validation score post-submission on \texttt{validation.csv}.
    \item \textbf{Routing vs. Coordinates}: Minor discrepancies between stated distance and coordinate haversine distance exist in raw data; distance was prioritized as the primary mileage signal.
    \item \textbf{Feature Parsimony}: \texttt{market\_index} and \texttt{quote\_signal} were omitted from the final feature set to ensure pipeline stability and avoid noisy residual estimates.
\end{itemize}

\section{Conclusion \& Verification Summary}
The complete pipeline has been verified through the official assessment scorer:
\begin{itemize}
    \item \textbf{Validation Predictions}: 12,000 positive, finite predictions generated and formatted as \texttt{load\_id,predicted\_rate}.
    \item \textbf{December Predictions}: 31 rows preserving original input specifications and dates.
    \item \textbf{Official Scorer}: \texttt{score.py} executes successfully with exit code 0, verifying schema compliance and outputting \texttt{candidate\_december.png}.
\end{itemize}

\end{document}